\section{The argument under review}
\label{sec:argument}

We first state the argument of \N\ in a form that separates what it observes from what it infers. Section numbers refer to \N.
The reconstruction was made twice, by us and independently by cross-check agent~H (Section~\ref{sec:agents}), with the same result.

\begin{notebox}
\textbf{Observations about the code}
\begin{enumerate}[label=(A\arabic*),leftmargin=3.2em]
\item The final theorem is a closed term of Mathlib's \lean{RiemannHypothesis}; the source has no \lean{axiom}, \lean{sorry}, \lean{native_decide}, \lean{unsafe} or \lean{opaque} declaration. \hfill (Abstract, \S12)
\item The proof term reaches it along the route \lean{riemannHypothesis} $\leftarrow$ \lean{rhOfAll} $\leftarrow$ \lean{allChi} $\leftarrow$ \lean{everyCharacter} $\leftarrow\cdots\leftarrow$
\lean{every_character_of_green_v_win} $\leftarrow$ \lean{zero_free_all_of_green_v_win} $\leftarrow$ \lean{zeroFreeAll_of_topG}. The theorem \lean{plugAll} is not on it. \hfill (\S1, \S7, \S11)
\item On this route the witness of the hypothesis bundle \lean{HypsAlt} is built with \texttt{geoCl anyGeometry}, where \lean{anyGeometry} is the constant \lean{True}; two further ``shape'' inputs are supplied by \lean{trivial}. \hfill (\S2, \S7)
\item The geometric datum is assembled from a degenerate model; the lattice conditions on it hold by \lean{rfl} because the seed is defined as a copy of the row-side data; the map from the lattice to sections is zero. \hfill (\S3--\S5)
\item The lemmas that would extract geometric content are never called; one geometric field is vacuously true for the model; the degree equation is never proved. \hfill (\S6)
\item The height of the divisor datum is zero when the datum is built and is then overwritten by an explicit formula \lean{hExpl6}, so ``height $=$ formula'' holds by \lean{rfl};
two inequalities about that height are then proved by \lean{divClausesR_of} from about forty matrix estimates. \hfill (\S7)
\end{enumerate}
\textbf{Interpretation and requirements}
\begin{enumerate}[label=(B\arabic*),leftmargin=3.2em]
\item Hence the divisor side and the row side are one object, and the inequality that should trap a geometric height between a pairing and a main term compares a quantity with itself and holds by reflexivity. \hfill (Abstract, \S8)
\item Hence what the kernel certifies is only the tautology ``RiemannHypothesis $=$ RiemannHypothesis'' under the supplied definitions. \hfill (Abstract, \S8)
\item The argument that the source \emph{names} needs an independent geometric height. A formula that was written in to fit is not one, and a geometric argument cannot be both essential and trivially discharged. \hfill (\S8, \S9)
\item The algebraic zero-free step is downstream of the bundle. Remove the bundle and the proof term has no route to the conclusion; so the objection that a kernel-checked closed term is a proof does not apply. \hfill (\S10, App.~A)
\end{enumerate}
\textbf{Conclusion}
\begin{enumerate}[label=(C),leftmargin=3.2em]
\item The release does not contain a proof of the Riemann Hypothesis. \hfill (Abstract, \S9, \S13)
\end{enumerate}
\end{notebox}

Two features of this argument should be noted at once.
First, (A1) is conceded by \N: the dispute is not about axioms, \lean{sorry} or the identity of the statement.
Second, \N\ anticipates the obvious reply, namely that a kernel-checked closed term is a derivation of its statement whatever its intermediate steps look like, and answers it with (B4).
The weight of the paper therefore rests on (B1)--(B4). Section~\ref{sec:facts} checks (A1)--(A6); Section~\ref{sec:inference} examines (B1)--(B4).
